\documentclass[11pt]{article}

    \usepackage[breakable]{tcolorbox}
    \usepackage{parskip} % Stop auto-indenting (to mimic markdown behaviour)
    

    % Basic figure setup, for now with no caption control since it's done
    % automatically by Pandoc (which extracts ![](path) syntax from Markdown).
    \usepackage{graphicx}
    % Maintain compatibility with old templates. Remove in nbconvert 6.0
    \let\Oldincludegraphics\includegraphics
    % Ensure that by default, figures have no caption (until we provide a
    % proper Figure object with a Caption API and a way to capture that
    % in the conversion process - todo).
    \usepackage{caption}
    \DeclareCaptionFormat{nocaption}{}
    \captionsetup{format=nocaption,aboveskip=0pt,belowskip=0pt}

    \usepackage{float}
    \floatplacement{figure}{H} % forces figures to be placed at the correct location
    \usepackage{xcolor} % Allow colors to be defined
    \usepackage{enumerate} % Needed for markdown enumerations to work
    \usepackage{geometry} % Used to adjust the document margins
    \usepackage{amsmath} % Equations
    \usepackage{amssymb} % Equations
    \usepackage{textcomp} % defines textquotesingle
    % Hack from http://tex.stackexchange.com/a/47451/13684:
    \AtBeginDocument{%
        \def\PYZsq{\textquotesingle}% Upright quotes in Pygmentized code
    }
    \usepackage{upquote} % Upright quotes for verbatim code
    \usepackage{eurosym} % defines \euro

    \usepackage{iftex}
    \ifPDFTeX
        \usepackage[T1]{fontenc}
        \IfFileExists{alphabeta.sty}{
              \usepackage{alphabeta}
          }{
              \usepackage[mathletters]{ucs}
              \usepackage[utf8x]{inputenc}
          }
    \else
        \usepackage{fontspec}
        \usepackage{unicode-math}
    \fi

    \usepackage{fancyvrb} % verbatim replacement that allows latex
    \usepackage[Export]{adjustbox} % Used to constrain images to a maximum size
    \adjustboxset{max size={0.9\linewidth}{0.9\paperheight}}

    % The hyperref package gives us a pdf with properly built
    % internal navigation ('pdf bookmarks' for the table of contents,
    % internal cross-reference links, web links for URLs, etc.)
    \usepackage{hyperref}
    % The default LaTeX title has an obnoxious amount of whitespace. By default,
    % titling removes some of it. It also provides customization options.
    \usepackage{titling}
    \usepackage{longtable} % longtable support required by pandoc >1.10
    \usepackage{booktabs}  % table support for pandoc > 1.12.2
    \usepackage{array}     % table support for pandoc >= 2.11.3
    \usepackage{calc}      % table minipage width calculation for pandoc >= 2.11.1
    \usepackage[inline]{enumitem} % IRkernel/repr support (it uses the enumerate* environment)
    \usepackage[normalem]{ulem} % ulem is needed to support strikethroughs (\sout)
                                % normalem makes italics be italics, not underlines
    \usepackage{mathrsfs}
    

    
    % Colors for the hyperref package
    \definecolor{urlcolor}{rgb}{0,.145,.698}
    \definecolor{linkcolor}{rgb}{.71,0.21,0.01}
    \definecolor{citecolor}{rgb}{.12,.54,.11}

    % ANSI colors
    \definecolor{ansi-black}{HTML}{3E424D}
    \definecolor{ansi-black-intense}{HTML}{282C36}
    \definecolor{ansi-red}{HTML}{E75C58}
    \definecolor{ansi-red-intense}{HTML}{B22B31}
    \definecolor{ansi-green}{HTML}{00A250}
    \definecolor{ansi-green-intense}{HTML}{007427}
    \definecolor{ansi-yellow}{HTML}{DDB62B}
    \definecolor{ansi-yellow-intense}{HTML}{B27D12}
    \definecolor{ansi-blue}{HTML}{208FFB}
    \definecolor{ansi-blue-intense}{HTML}{0065CA}
    \definecolor{ansi-magenta}{HTML}{D160C4}
    \definecolor{ansi-magenta-intense}{HTML}{A03196}
    \definecolor{ansi-cyan}{HTML}{60C6C8}
    \definecolor{ansi-cyan-intense}{HTML}{258F8F}
    \definecolor{ansi-white}{HTML}{C5C1B4}
    \definecolor{ansi-white-intense}{HTML}{A1A6B2}
    \definecolor{ansi-default-inverse-fg}{HTML}{FFFFFF}
    \definecolor{ansi-default-inverse-bg}{HTML}{000000}

    % common color for the border for error outputs.
    \definecolor{outerrorbackground}{HTML}{FFDFDF}

    % commands and environments needed by pandoc snippets
    % extracted from the output of `pandoc -s`
    \providecommand{\tightlist}{%
      \setlength{\itemsep}{0pt}\setlength{\parskip}{0pt}}
    \DefineVerbatimEnvironment{Highlighting}{Verbatim}{commandchars=\\\{\}}
    % Add ',fontsize=\small' for more characters per line
    \newenvironment{Shaded}{}{}
    \newcommand{\KeywordTok}[1]{\textcolor[rgb]{0.00,0.44,0.13}{\textbf{{#1}}}}
    \newcommand{\DataTypeTok}[1]{\textcolor[rgb]{0.56,0.13,0.00}{{#1}}}
    \newcommand{\DecValTok}[1]{\textcolor[rgb]{0.25,0.63,0.44}{{#1}}}
    \newcommand{\BaseNTok}[1]{\textcolor[rgb]{0.25,0.63,0.44}{{#1}}}
    \newcommand{\FloatTok}[1]{\textcolor[rgb]{0.25,0.63,0.44}{{#1}}}
    \newcommand{\CharTok}[1]{\textcolor[rgb]{0.25,0.44,0.63}{{#1}}}
    \newcommand{\StringTok}[1]{\textcolor[rgb]{0.25,0.44,0.63}{{#1}}}
    \newcommand{\CommentTok}[1]{\textcolor[rgb]{0.38,0.63,0.69}{\textit{{#1}}}}
    \newcommand{\OtherTok}[1]{\textcolor[rgb]{0.00,0.44,0.13}{{#1}}}
    \newcommand{\AlertTok}[1]{\textcolor[rgb]{1.00,0.00,0.00}{\textbf{{#1}}}}
    \newcommand{\FunctionTok}[1]{\textcolor[rgb]{0.02,0.16,0.49}{{#1}}}
    \newcommand{\RegionMarkerTok}[1]{{#1}}
    \newcommand{\ErrorTok}[1]{\textcolor[rgb]{1.00,0.00,0.00}{\textbf{{#1}}}}
    \newcommand{\NormalTok}[1]{{#1}}

    % Additional commands for more recent versions of Pandoc
    \newcommand{\ConstantTok}[1]{\textcolor[rgb]{0.53,0.00,0.00}{{#1}}}
    \newcommand{\SpecialCharTok}[1]{\textcolor[rgb]{0.25,0.44,0.63}{{#1}}}
    \newcommand{\VerbatimStringTok}[1]{\textcolor[rgb]{0.25,0.44,0.63}{{#1}}}
    \newcommand{\SpecialStringTok}[1]{\textcolor[rgb]{0.73,0.40,0.53}{{#1}}}
    \newcommand{\ImportTok}[1]{{#1}}
    \newcommand{\DocumentationTok}[1]{\textcolor[rgb]{0.73,0.13,0.13}{\textit{{#1}}}}
    \newcommand{\AnnotationTok}[1]{\textcolor[rgb]{0.38,0.63,0.69}{\textbf{\textit{{#1}}}}}
    \newcommand{\CommentVarTok}[1]{\textcolor[rgb]{0.38,0.63,0.69}{\textbf{\textit{{#1}}}}}
    \newcommand{\VariableTok}[1]{\textcolor[rgb]{0.10,0.09,0.49}{{#1}}}
    \newcommand{\ControlFlowTok}[1]{\textcolor[rgb]{0.00,0.44,0.13}{\textbf{{#1}}}}
    \newcommand{\OperatorTok}[1]{\textcolor[rgb]{0.40,0.40,0.40}{{#1}}}
    \newcommand{\BuiltInTok}[1]{{#1}}
    \newcommand{\ExtensionTok}[1]{{#1}}
    \newcommand{\PreprocessorTok}[1]{\textcolor[rgb]{0.74,0.48,0.00}{{#1}}}
    \newcommand{\AttributeTok}[1]{\textcolor[rgb]{0.49,0.56,0.16}{{#1}}}
    \newcommand{\InformationTok}[1]{\textcolor[rgb]{0.38,0.63,0.69}{\textbf{\textit{{#1}}}}}
    \newcommand{\WarningTok}[1]{\textcolor[rgb]{0.38,0.63,0.69}{\textbf{\textit{{#1}}}}}


    % Define a nice break command that doesn't care if a line doesn't already
    % exist.
    \def\br{\hspace*{\fill} \\* }
    % Math Jax compatibility definitions
    \def\gt{>}
    \def\lt{<}
    \let\Oldtex\TeX
    \let\Oldlatex\LaTeX
    \renewcommand{\TeX}{\textrm{\Oldtex}}
    \renewcommand{\LaTeX}{\textrm{\Oldlatex}}
    % Document parameters
    % Document title
    \title{Copia\_de\_Copia\_de\_Tuto\_1\_2\_Open\_Data\_access\_with\_GWpy}
    
    
    
    
    
% Pygments definitions
\makeatletter
\def\PY@reset{\let\PY@it=\relax \let\PY@bf=\relax%
    \let\PY@ul=\relax \let\PY@tc=\relax%
    \let\PY@bc=\relax \let\PY@ff=\relax}
\def\PY@tok#1{\csname PY@tok@#1\endcsname}
\def\PY@toks#1+{\ifx\relax#1\empty\else%
    \PY@tok{#1}\expandafter\PY@toks\fi}
\def\PY@do#1{\PY@bc{\PY@tc{\PY@ul{%
    \PY@it{\PY@bf{\PY@ff{#1}}}}}}}
\def\PY#1#2{\PY@reset\PY@toks#1+\relax+\PY@do{#2}}

\@namedef{PY@tok@w}{\def\PY@tc##1{\textcolor[rgb]{0.73,0.73,0.73}{##1}}}
\@namedef{PY@tok@c}{\let\PY@it=\textit\def\PY@tc##1{\textcolor[rgb]{0.24,0.48,0.48}{##1}}}
\@namedef{PY@tok@cp}{\def\PY@tc##1{\textcolor[rgb]{0.61,0.40,0.00}{##1}}}
\@namedef{PY@tok@k}{\let\PY@bf=\textbf\def\PY@tc##1{\textcolor[rgb]{0.00,0.50,0.00}{##1}}}
\@namedef{PY@tok@kp}{\def\PY@tc##1{\textcolor[rgb]{0.00,0.50,0.00}{##1}}}
\@namedef{PY@tok@kt}{\def\PY@tc##1{\textcolor[rgb]{0.69,0.00,0.25}{##1}}}
\@namedef{PY@tok@o}{\def\PY@tc##1{\textcolor[rgb]{0.40,0.40,0.40}{##1}}}
\@namedef{PY@tok@ow}{\let\PY@bf=\textbf\def\PY@tc##1{\textcolor[rgb]{0.67,0.13,1.00}{##1}}}
\@namedef{PY@tok@nb}{\def\PY@tc##1{\textcolor[rgb]{0.00,0.50,0.00}{##1}}}
\@namedef{PY@tok@nf}{\def\PY@tc##1{\textcolor[rgb]{0.00,0.00,1.00}{##1}}}
\@namedef{PY@tok@nc}{\let\PY@bf=\textbf\def\PY@tc##1{\textcolor[rgb]{0.00,0.00,1.00}{##1}}}
\@namedef{PY@tok@nn}{\let\PY@bf=\textbf\def\PY@tc##1{\textcolor[rgb]{0.00,0.00,1.00}{##1}}}
\@namedef{PY@tok@ne}{\let\PY@bf=\textbf\def\PY@tc##1{\textcolor[rgb]{0.80,0.25,0.22}{##1}}}
\@namedef{PY@tok@nv}{\def\PY@tc##1{\textcolor[rgb]{0.10,0.09,0.49}{##1}}}
\@namedef{PY@tok@no}{\def\PY@tc##1{\textcolor[rgb]{0.53,0.00,0.00}{##1}}}
\@namedef{PY@tok@nl}{\def\PY@tc##1{\textcolor[rgb]{0.46,0.46,0.00}{##1}}}
\@namedef{PY@tok@ni}{\let\PY@bf=\textbf\def\PY@tc##1{\textcolor[rgb]{0.44,0.44,0.44}{##1}}}
\@namedef{PY@tok@na}{\def\PY@tc##1{\textcolor[rgb]{0.41,0.47,0.13}{##1}}}
\@namedef{PY@tok@nt}{\let\PY@bf=\textbf\def\PY@tc##1{\textcolor[rgb]{0.00,0.50,0.00}{##1}}}
\@namedef{PY@tok@nd}{\def\PY@tc##1{\textcolor[rgb]{0.67,0.13,1.00}{##1}}}
\@namedef{PY@tok@s}{\def\PY@tc##1{\textcolor[rgb]{0.73,0.13,0.13}{##1}}}
\@namedef{PY@tok@sd}{\let\PY@it=\textit\def\PY@tc##1{\textcolor[rgb]{0.73,0.13,0.13}{##1}}}
\@namedef{PY@tok@si}{\let\PY@bf=\textbf\def\PY@tc##1{\textcolor[rgb]{0.64,0.35,0.47}{##1}}}
\@namedef{PY@tok@se}{\let\PY@bf=\textbf\def\PY@tc##1{\textcolor[rgb]{0.67,0.36,0.12}{##1}}}
\@namedef{PY@tok@sr}{\def\PY@tc##1{\textcolor[rgb]{0.64,0.35,0.47}{##1}}}
\@namedef{PY@tok@ss}{\def\PY@tc##1{\textcolor[rgb]{0.10,0.09,0.49}{##1}}}
\@namedef{PY@tok@sx}{\def\PY@tc##1{\textcolor[rgb]{0.00,0.50,0.00}{##1}}}
\@namedef{PY@tok@m}{\def\PY@tc##1{\textcolor[rgb]{0.40,0.40,0.40}{##1}}}
\@namedef{PY@tok@gh}{\let\PY@bf=\textbf\def\PY@tc##1{\textcolor[rgb]{0.00,0.00,0.50}{##1}}}
\@namedef{PY@tok@gu}{\let\PY@bf=\textbf\def\PY@tc##1{\textcolor[rgb]{0.50,0.00,0.50}{##1}}}
\@namedef{PY@tok@gd}{\def\PY@tc##1{\textcolor[rgb]{0.63,0.00,0.00}{##1}}}
\@namedef{PY@tok@gi}{\def\PY@tc##1{\textcolor[rgb]{0.00,0.52,0.00}{##1}}}
\@namedef{PY@tok@gr}{\def\PY@tc##1{\textcolor[rgb]{0.89,0.00,0.00}{##1}}}
\@namedef{PY@tok@ge}{\let\PY@it=\textit}
\@namedef{PY@tok@gs}{\let\PY@bf=\textbf}
\@namedef{PY@tok@gp}{\let\PY@bf=\textbf\def\PY@tc##1{\textcolor[rgb]{0.00,0.00,0.50}{##1}}}
\@namedef{PY@tok@go}{\def\PY@tc##1{\textcolor[rgb]{0.44,0.44,0.44}{##1}}}
\@namedef{PY@tok@gt}{\def\PY@tc##1{\textcolor[rgb]{0.00,0.27,0.87}{##1}}}
\@namedef{PY@tok@err}{\def\PY@bc##1{{\setlength{\fboxsep}{\string -\fboxrule}\fcolorbox[rgb]{1.00,0.00,0.00}{1,1,1}{\strut ##1}}}}
\@namedef{PY@tok@kc}{\let\PY@bf=\textbf\def\PY@tc##1{\textcolor[rgb]{0.00,0.50,0.00}{##1}}}
\@namedef{PY@tok@kd}{\let\PY@bf=\textbf\def\PY@tc##1{\textcolor[rgb]{0.00,0.50,0.00}{##1}}}
\@namedef{PY@tok@kn}{\let\PY@bf=\textbf\def\PY@tc##1{\textcolor[rgb]{0.00,0.50,0.00}{##1}}}
\@namedef{PY@tok@kr}{\let\PY@bf=\textbf\def\PY@tc##1{\textcolor[rgb]{0.00,0.50,0.00}{##1}}}
\@namedef{PY@tok@bp}{\def\PY@tc##1{\textcolor[rgb]{0.00,0.50,0.00}{##1}}}
\@namedef{PY@tok@fm}{\def\PY@tc##1{\textcolor[rgb]{0.00,0.00,1.00}{##1}}}
\@namedef{PY@tok@vc}{\def\PY@tc##1{\textcolor[rgb]{0.10,0.09,0.49}{##1}}}
\@namedef{PY@tok@vg}{\def\PY@tc##1{\textcolor[rgb]{0.10,0.09,0.49}{##1}}}
\@namedef{PY@tok@vi}{\def\PY@tc##1{\textcolor[rgb]{0.10,0.09,0.49}{##1}}}
\@namedef{PY@tok@vm}{\def\PY@tc##1{\textcolor[rgb]{0.10,0.09,0.49}{##1}}}
\@namedef{PY@tok@sa}{\def\PY@tc##1{\textcolor[rgb]{0.73,0.13,0.13}{##1}}}
\@namedef{PY@tok@sb}{\def\PY@tc##1{\textcolor[rgb]{0.73,0.13,0.13}{##1}}}
\@namedef{PY@tok@sc}{\def\PY@tc##1{\textcolor[rgb]{0.73,0.13,0.13}{##1}}}
\@namedef{PY@tok@dl}{\def\PY@tc##1{\textcolor[rgb]{0.73,0.13,0.13}{##1}}}
\@namedef{PY@tok@s2}{\def\PY@tc##1{\textcolor[rgb]{0.73,0.13,0.13}{##1}}}
\@namedef{PY@tok@sh}{\def\PY@tc##1{\textcolor[rgb]{0.73,0.13,0.13}{##1}}}
\@namedef{PY@tok@s1}{\def\PY@tc##1{\textcolor[rgb]{0.73,0.13,0.13}{##1}}}
\@namedef{PY@tok@mb}{\def\PY@tc##1{\textcolor[rgb]{0.40,0.40,0.40}{##1}}}
\@namedef{PY@tok@mf}{\def\PY@tc##1{\textcolor[rgb]{0.40,0.40,0.40}{##1}}}
\@namedef{PY@tok@mh}{\def\PY@tc##1{\textcolor[rgb]{0.40,0.40,0.40}{##1}}}
\@namedef{PY@tok@mi}{\def\PY@tc##1{\textcolor[rgb]{0.40,0.40,0.40}{##1}}}
\@namedef{PY@tok@il}{\def\PY@tc##1{\textcolor[rgb]{0.40,0.40,0.40}{##1}}}
\@namedef{PY@tok@mo}{\def\PY@tc##1{\textcolor[rgb]{0.40,0.40,0.40}{##1}}}
\@namedef{PY@tok@ch}{\let\PY@it=\textit\def\PY@tc##1{\textcolor[rgb]{0.24,0.48,0.48}{##1}}}
\@namedef{PY@tok@cm}{\let\PY@it=\textit\def\PY@tc##1{\textcolor[rgb]{0.24,0.48,0.48}{##1}}}
\@namedef{PY@tok@cpf}{\let\PY@it=\textit\def\PY@tc##1{\textcolor[rgb]{0.24,0.48,0.48}{##1}}}
\@namedef{PY@tok@c1}{\let\PY@it=\textit\def\PY@tc##1{\textcolor[rgb]{0.24,0.48,0.48}{##1}}}
\@namedef{PY@tok@cs}{\let\PY@it=\textit\def\PY@tc##1{\textcolor[rgb]{0.24,0.48,0.48}{##1}}}

\def\PYZbs{\char`\\}
\def\PYZus{\char`\_}
\def\PYZob{\char`\{}
\def\PYZcb{\char`\}}
\def\PYZca{\char`\^}
\def\PYZam{\char`\&}
\def\PYZlt{\char`\<}
\def\PYZgt{\char`\>}
\def\PYZsh{\char`\#}
\def\PYZpc{\char`\%}
\def\PYZdl{\char`\$}
\def\PYZhy{\char`\-}
\def\PYZsq{\char`\'}
\def\PYZdq{\char`\"}
\def\PYZti{\char`\~}
% for compatibility with earlier versions
\def\PYZat{@}
\def\PYZlb{[}
\def\PYZrb{]}
\makeatother


    % For linebreaks inside Verbatim environment from package fancyvrb.
    \makeatletter
        \newbox\Wrappedcontinuationbox
        \newbox\Wrappedvisiblespacebox
        \newcommand*\Wrappedvisiblespace {\textcolor{red}{\textvisiblespace}}
        \newcommand*\Wrappedcontinuationsymbol {\textcolor{red}{\llap{\tiny$\m@th\hookrightarrow$}}}
        \newcommand*\Wrappedcontinuationindent {3ex }
        \newcommand*\Wrappedafterbreak {\kern\Wrappedcontinuationindent\copy\Wrappedcontinuationbox}
        % Take advantage of the already applied Pygments mark-up to insert
        % potential linebreaks for TeX processing.
        %        {, <, #, %, $, ' and ": go to next line.
        %        _, }, ^, &, >, - and ~: stay at end of broken line.
        % Use of \textquotesingle for straight quote.
        \newcommand*\Wrappedbreaksatspecials {%
            \def\PYGZus{\discretionary{\char`\_}{\Wrappedafterbreak}{\char`\_}}%
            \def\PYGZob{\discretionary{}{\Wrappedafterbreak\char`\{}{\char`\{}}%
            \def\PYGZcb{\discretionary{\char`\}}{\Wrappedafterbreak}{\char`\}}}%
            \def\PYGZca{\discretionary{\char`\^}{\Wrappedafterbreak}{\char`\^}}%
            \def\PYGZam{\discretionary{\char`\&}{\Wrappedafterbreak}{\char`\&}}%
            \def\PYGZlt{\discretionary{}{\Wrappedafterbreak\char`\<}{\char`\<}}%
            \def\PYGZgt{\discretionary{\char`\>}{\Wrappedafterbreak}{\char`\>}}%
            \def\PYGZsh{\discretionary{}{\Wrappedafterbreak\char`\#}{\char`\#}}%
            \def\PYGZpc{\discretionary{}{\Wrappedafterbreak\char`\%}{\char`\%}}%
            \def\PYGZdl{\discretionary{}{\Wrappedafterbreak\char`\$}{\char`\$}}%
            \def\PYGZhy{\discretionary{\char`\-}{\Wrappedafterbreak}{\char`\-}}%
            \def\PYGZsq{\discretionary{}{\Wrappedafterbreak\textquotesingle}{\textquotesingle}}%
            \def\PYGZdq{\discretionary{}{\Wrappedafterbreak\char`\"}{\char`\"}}%
            \def\PYGZti{\discretionary{\char`\~}{\Wrappedafterbreak}{\char`\~}}%
        }
        % Some characters . , ; ? ! / are not pygmentized.
        % This macro makes them "active" and they will insert potential linebreaks
        \newcommand*\Wrappedbreaksatpunct {%
            \lccode`\~`\.\lowercase{\def~}{\discretionary{\hbox{\char`\.}}{\Wrappedafterbreak}{\hbox{\char`\.}}}%
            \lccode`\~`\,\lowercase{\def~}{\discretionary{\hbox{\char`\,}}{\Wrappedafterbreak}{\hbox{\char`\,}}}%
            \lccode`\~`\;\lowercase{\def~}{\discretionary{\hbox{\char`\;}}{\Wrappedafterbreak}{\hbox{\char`\;}}}%
            \lccode`\~`\:\lowercase{\def~}{\discretionary{\hbox{\char`\:}}{\Wrappedafterbreak}{\hbox{\char`\:}}}%
            \lccode`\~`\?\lowercase{\def~}{\discretionary{\hbox{\char`\?}}{\Wrappedafterbreak}{\hbox{\char`\?}}}%
            \lccode`\~`\!\lowercase{\def~}{\discretionary{\hbox{\char`\!}}{\Wrappedafterbreak}{\hbox{\char`\!}}}%
            \lccode`\~`\/\lowercase{\def~}{\discretionary{\hbox{\char`\/}}{\Wrappedafterbreak}{\hbox{\char`\/}}}%
            \catcode`\.\active
            \catcode`\,\active
            \catcode`\;\active
            \catcode`\:\active
            \catcode`\?\active
            \catcode`\!\active
            \catcode`\/\active
            \lccode`\~`\~
        }
    \makeatother

    \let\OriginalVerbatim=\Verbatim
    \makeatletter
    \renewcommand{\Verbatim}[1][1]{%
        %\parskip\z@skip
        \sbox\Wrappedcontinuationbox {\Wrappedcontinuationsymbol}%
        \sbox\Wrappedvisiblespacebox {\FV@SetupFont\Wrappedvisiblespace}%
        \def\FancyVerbFormatLine ##1{\hsize\linewidth
            \vtop{\raggedright\hyphenpenalty\z@\exhyphenpenalty\z@
                \doublehyphendemerits\z@\finalhyphendemerits\z@
                \strut ##1\strut}%
        }%
        % If the linebreak is at a space, the latter will be displayed as visible
        % space at end of first line, and a continuation symbol starts next line.
        % Stretch/shrink are however usually zero for typewriter font.
        \def\FV@Space {%
            \nobreak\hskip\z@ plus\fontdimen3\font minus\fontdimen4\font
            \discretionary{\copy\Wrappedvisiblespacebox}{\Wrappedafterbreak}
            {\kern\fontdimen2\font}%
        }%

        % Allow breaks at special characters using \PYG... macros.
        \Wrappedbreaksatspecials
        % Breaks at punctuation characters . , ; ? ! and / need catcode=\active
        \OriginalVerbatim[#1,codes*=\Wrappedbreaksatpunct]%
    }
    \makeatother

    % Exact colors from NB
    \definecolor{incolor}{HTML}{303F9F}
    \definecolor{outcolor}{HTML}{D84315}
    \definecolor{cellborder}{HTML}{CFCFCF}
    \definecolor{cellbackground}{HTML}{F7F7F7}

    % prompt
    \makeatletter
    \newcommand{\boxspacing}{\kern\kvtcb@left@rule\kern\kvtcb@boxsep}
    \makeatother
    \newcommand{\prompt}[4]{
        {\ttfamily\llap{{\color{#2}[#3]:\hspace{3pt}#4}}\vspace{-\baselineskip}}
    }
    

    
    % Prevent overflowing lines due to hard-to-break entities
    \sloppy
    % Setup hyperref package
    \hypersetup{
      breaklinks=true,  % so long urls are correctly broken across lines
      colorlinks=true,
      urlcolor=urlcolor,
      linkcolor=linkcolor,
      citecolor=citecolor,
      }
    % Slightly bigger margins than the latex defaults
    
    \geometry{verbose,tmargin=1in,bmargin=1in,lmargin=1in,rmargin=1in}
    
    

\begin{document}
    
    \maketitle
    
    

    
    \hypertarget{gravitational-wave-open-data-workshop-4}{%
\section{Gravitational Wave Open Data Workshop
\#4}\label{gravitational-wave-open-data-workshop-4}}

\hypertarget{tutorial-1.2-introduction-to-gwpy}{%
\paragraph{Tutorial 1.2: Introduction to
GWpy}\label{tutorial-1.2-introduction-to-gwpy}}

This tutorial will briefly describe GWpy, a python package for
gravitational astrophysics, and walk-through how you can use this to
speed up access to, and processing of, GWOSC data.

\href{https://colab.research.google.com/github/gw-odw/odw-2021/blob/master/Day_1/Tuto\%201.2\%20Open\%20Data\%20access\%20with\%20GWpy.ipynb}{Click
this link to view this tutorial in Google Colaboratory}

This notebook was generated using python 3.8, but should work on python
2.7, 3.6, or 3.7.

    \hypertarget{installation-execute-only-if-running-on-a-cloud-platform-like-google-colab-or-if-you-havent-done-the-installation-already}{%
\subsection{Installation (execute only if running on a cloud platform,
like Google Colab, or if you haven't done the installation
already!)}\label{installation-execute-only-if-running-on-a-cloud-platform-like-google-colab-or-if-you-havent-done-the-installation-already}}

Note: we use
\href{https://docs.python.org/3.6/installing/}{\texttt{pip}}, but
\textbf{it is recommended} to use
\href{https://docs.ligo.org/lscsoft/conda/}{conda} on your own machine,
as explained in the
\href{https://github.com/gw-odw/odw-2021\%7C/blob/master/setup.md}{installation
instructions}. This usage might look a little different than normal,
simply because we want to do this directly from the notebook.

    \begin{tcolorbox}[breakable, size=fbox, boxrule=1pt, pad at break*=1mm,colback=cellbackground, colframe=cellborder]
\prompt{In}{incolor}{2}{\boxspacing}
\begin{Verbatim}[commandchars=\\\{\}]
\PY{c+c1}{\PYZsh{} \PYZhy{}\PYZhy{} Uncomment following line if running in Google Colab}
\PY{o}{!}\PY{+w}{ }pip\PY{+w}{ }install\PY{+w}{ }\PYZhy{}q\PY{+w}{ }\PY{l+s+s1}{\PYZsq{}gwpy==3.0.12\PYZsq{}}
\end{Verbatim}
\end{tcolorbox}

    \begin{Verbatim}[commandchars=\\\{\}]
   \textcolor{ansi-black-intense}{━━━━━━━━━━━━━━━━━━━━━━━━━━━━━━━━━━━━━━━━} \textcolor{ansi-green}{1.4/1.4 MB}
\textcolor{ansi-red}{11.0 MB/s} eta \textcolor{ansi-cyan}{0:00:00}
   \textcolor{ansi-black-intense}{━━━━━━━━━━━━━━━━━━━━━━━━━━━━━━━━━━━━━━━━} \textcolor{ansi-green}{315.5/315.5 kB}
\textcolor{ansi-red}{24.1 MB/s} eta \textcolor{ansi-cyan}{0:00:00}
   \textcolor{ansi-black-intense}{━━━━━━━━━━━━━━━━━━━━━━━━━━━━━━━━━━━━━━━━} \textcolor{ansi-green}{45.2/45.2 kB}
\textcolor{ansi-red}{3.4 MB/s} eta \textcolor{ansi-cyan}{0:00:00}
   \textcolor{ansi-black-intense}{━━━━━━━━━━━━━━━━━━━━━━━━━━━━━━━━━━━━━━━━} \textcolor{ansi-green}{131.1/131.1 kB}
\textcolor{ansi-red}{10.5 MB/s} eta \textcolor{ansi-cyan}{0:00:00}
   \textcolor{ansi-black-intense}{━━━━━━━━━━━━━━━━━━━━━━━━━━━━━━━━━━━━━━━━} \textcolor{ansi-green}{4.5/4.5 MB}
\textcolor{ansi-red}{101.3 MB/s} eta \textcolor{ansi-cyan}{0:00:00}
\textcolor{ansi-red}{ERROR: pip's dependency resolver does not currently take into account
all the packages that are installed. This behaviour is the source of the
following dependency conflicts.
pyopenssl 24.2.1 requires cryptography<44,>=41.0.5, but you have cryptography
46.0.3 which is incompatible.
pydrive2 1.21.3 requires cryptography<44, but you have cryptography 46.0.3 which
is incompatible.}\textcolor{ansi-red}{
}
    \end{Verbatim}

    \textbf{Important:} With Google Colab, you may need to restart the
runtime after running the cell above.

    \hypertarget{initialization}{%
\subsection{Initialization}\label{initialization}}

    \begin{tcolorbox}[breakable, size=fbox, boxrule=1pt, pad at break*=1mm,colback=cellbackground, colframe=cellborder]
\prompt{In}{incolor}{3}{\boxspacing}
\begin{Verbatim}[commandchars=\\\{\}]
\PY{k+kn}{import} \PY{n+nn}{gwpy}
\PY{n+nb}{print}\PY{p}{(}\PY{n}{gwpy}\PY{o}{.}\PY{n}{\PYZus{}\PYZus{}version\PYZus{}\PYZus{}}\PY{p}{)}

\PY{k+kn}{from} \PY{n+nn}{gwpy}\PY{n+nn}{.}\PY{n+nn}{timeseries} \PY{k+kn}{import} \PY{n}{TimeSeries}
\PY{n+nb}{print}\PY{p}{(}\PY{l+s+s2}{\PYZdq{}}\PY{l+s+s2}{gwpy importado OK:}\PY{l+s+s2}{\PYZdq{}}\PY{p}{,} \PY{n}{TimeSeries}\PY{p}{)}
\end{Verbatim}
\end{tcolorbox}

    \begin{Verbatim}[commandchars=\\\{\}]
3.0.12
gwpy importado OK: <class 'gwpy.timeseries.timeseries.TimeSeries'>
    \end{Verbatim}

    \hypertarget{a-note-on-object-oriented-programming}{%
\subsection{A note on object-oriented
programming}\label{a-note-on-object-oriented-programming}}

Before we dive too deeply, it's worth a quick aside on object-oriented
programming (OOP).
\href{https://gwpy.github.io/docs/stable/index.html}{GWpy} is heavily
object-oriented, meaning almost all of the code you run using GWpy is
based around an object of some type, e.g.~\texttt{TimeSeries}. Most of
the methods (functions) we will use are attached to an object, rather
than standing alone, meaning you should have a pretty good idea of what
sort of data you are dealing with (without having to read the
documentation!).

For a quick overview of object-oriented programming in Python, see
\href{https://jeffknupp.com/blog/2014/06/18/improve-your-python-python-classes-and-object-oriented-programming/}{this
blog post by Jeff Knupp}.

    \hypertarget{handling-data-in-the-time-domain}{%
\subsection{Handling data in the time
domain}\label{handling-data-in-the-time-domain}}

    \hypertarget{finding-open-data}{%
\paragraph{Finding open data}\label{finding-open-data}}

We have seen already that the \texttt{gwosc} module can be used to query
for what data are available on
\href{https://www.gw-openscience.org/data/}{GWOSC}. The next thing to do
is to actually read some open data. Let's try to get some for GW190412,
the first detection of a gravitational-wave signal from an unequal-mass
BBH (binary black hole system).

We can use the
\href{https://gwpy.github.io/docs/stable/api/gwpy.timeseries.TimeSeries.html\#gwpy.timeseries.TimeSeries.fetch_open_data}{\texttt{TimeSeries.fetch\_open\_data}}
method to download data directly from https://www.gw-openscience.org,
but we need to know the GPS times. We can query for the GPS time of an
event as follows:

    \begin{tcolorbox}[breakable, size=fbox, boxrule=1pt, pad at break*=1mm,colback=cellbackground, colframe=cellborder]
\prompt{In}{incolor}{4}{\boxspacing}
\begin{Verbatim}[commandchars=\\\{\}]
\PY{k+kn}{from} \PY{n+nn}{gwosc}\PY{n+nn}{.}\PY{n+nn}{datasets} \PY{k+kn}{import} \PY{n}{event\PYZus{}gps}
\PY{n}{gps} \PY{o}{=} \PY{n}{event\PYZus{}gps}\PY{p}{(}\PY{l+s+s1}{\PYZsq{}}\PY{l+s+s1}{GW190412}\PY{l+s+s1}{\PYZsq{}}\PY{p}{)}
\PY{n+nb}{print}\PY{p}{(}\PY{n}{gps}\PY{p}{)}
\end{Verbatim}
\end{tcolorbox}

    \begin{Verbatim}[commandchars=\\\{\}]
1239082262.1
    \end{Verbatim}

    Now we can build a \texttt{{[}start,\ end)} GPS segment to 10 seconds
around this time, using integers for convenience:

    \begin{tcolorbox}[breakable, size=fbox, boxrule=1pt, pad at break*=1mm,colback=cellbackground, colframe=cellborder]
\prompt{In}{incolor}{5}{\boxspacing}
\begin{Verbatim}[commandchars=\\\{\}]
\PY{n}{segment} \PY{o}{=} \PY{p}{(}\PY{n+nb}{int}\PY{p}{(}\PY{n}{gps}\PY{p}{)}\PY{o}{\PYZhy{}}\PY{l+m+mi}{5}\PY{p}{,} \PY{n+nb}{int}\PY{p}{(}\PY{n}{gps}\PY{p}{)}\PY{o}{+}\PY{l+m+mi}{5}\PY{p}{)}
\PY{n+nb}{print}\PY{p}{(}\PY{n}{segment}\PY{p}{)}
\end{Verbatim}
\end{tcolorbox}

    \begin{Verbatim}[commandchars=\\\{\}]
(1239082257, 1239082267)
    \end{Verbatim}

    and can now query for the full data. For this example we choose to
retrieve data for the LIGO-Livingston interferometer, using the
identifier \texttt{\textquotesingle{}L1\textquotesingle{}}. We could
have chosen any of

\begin{itemize}
\tightlist
\item
  \texttt{\textquotesingle{}G1}' - GEO600
\item
  \texttt{\textquotesingle{}H1\textquotesingle{}} - LIGO-Hanford
\item
  \texttt{\textquotesingle{}L1\textquotesingle{}} - LIGO-Livingston
\item
  \texttt{\textquotesingle{}V1\textquotesingle{}} - (Advanced) Virgo
\end{itemize}

In the future, the Japanese observatory KAGRA will come online, with the
identifier \texttt{\textquotesingle{}K1\textquotesingle{}}.

    \begin{tcolorbox}[breakable, size=fbox, boxrule=1pt, pad at break*=1mm,colback=cellbackground, colframe=cellborder]
\prompt{In}{incolor}{6}{\boxspacing}
\begin{Verbatim}[commandchars=\\\{\}]
\PY{k+kn}{from} \PY{n+nn}{gwpy}\PY{n+nn}{.}\PY{n+nn}{timeseries} \PY{k+kn}{import} \PY{n}{TimeSeries}
\PY{n}{ldata} \PY{o}{=} \PY{n}{TimeSeries}\PY{o}{.}\PY{n}{fetch\PYZus{}open\PYZus{}data}\PY{p}{(}\PY{l+s+s1}{\PYZsq{}}\PY{l+s+s1}{L1}\PY{l+s+s1}{\PYZsq{}}\PY{p}{,} \PY{o}{*}\PY{n}{segment}\PY{p}{,} \PY{n}{verbose}\PY{o}{=}\PY{k+kc}{True}\PY{p}{)}
\PY{n+nb}{print}\PY{p}{(}\PY{n}{ldata}\PY{p}{)}
\end{Verbatim}
\end{tcolorbox}

    \begin{Verbatim}[commandchars=\\\{\}]
Fetched 1 URLs from gwosc.org for [1239082257 .. 1239082267))
Reading data{\ldots} [Done]
TimeSeries([-8.42597565e-19, -8.52437103e-19, -8.60738804e-19,
            {\ldots},  1.38850270e-19,  1.37760541e-19,
             1.38094202e-19]
           unit: dimensionless,
           t0: 1239082257.0 s,
           dt: 0.000244140625 s,
           name: Strain,
           channel: None)
    \end{Verbatim}

    \hypertarget{the-verbosetrue-flag-lets-us-see-that-gwpy-has-discovered-one-file-that-provides-the-data-for-the-given-interval-downloaded-it-and-loaded-the-data.}{%
\subparagraph{\texorpdfstring{The \texttt{verbose=True} flag lets us see
that GWpy has discovered one file that provides the data for the given
interval, downloaded it, and loaded the
data.}{The verbose=True flag lets us see that GWpy has discovered one file that provides the data for the given interval, downloaded it, and loaded the data.}}\label{the-verbosetrue-flag-lets-us-see-that-gwpy-has-discovered-one-file-that-provides-the-data-for-the-given-interval-downloaded-it-and-loaded-the-data.}}

The files are not stored permanently, so next time you do the same call,
it will be downloaded again, however, if you know you might repeat the
same call many times, you can use \texttt{cache=True} to store the file
on your computer.

Notes:

\begin{itemize}
\item
  To control the dataset from which your data come from you can use the
  `dataset' keyword. It is recommended to use data from a run if they
  are available, because they contain the most updated version of the
  calibration. To check which was the run at the time you want to query,
  it is enough to use the find\_datasets method of the package gwosc
  discussed in the previous tutorial, specifying
  (type=`run',segment=segment). For example, for the segment we are
  using in this tutorial, the two available datasets are
  {[}`O3a\_16KHZ\_R1', `O3a\_4KHZ\_R1'{]}. For the sampling at 4 kHz,
  the complete command to get data from this dataset is then:
  \_TimeSeries.fetch\_open\_data(`L1', *segment, verbose=True,
  dataset=`O3a\_4KHZ\_R1')\_
\item
  To read data from a local file instead of from the GWOSC server, we
  can use
  \href{https://gwpy.github.io/docs/stable/api/gwpy.timeseries.TimeSeries.html\#gwpy.timeseries.TimeSeries.read}{\texttt{TimeSeries.read}}
  method.
\end{itemize}

We have now downloaded real LIGO data for GW190412! These are the actual
data used in the analysis that discovered the first binary black hole
merger.

    To sanity check things, we can easily make a plot, using the
\href{https://gwpy.github.io/docs/stable/timeseries/plot.html}{\texttt{plot()}}
method of the \texttt{data} \texttt{TimeSeries}.

    Since this is the first time we are plotting something in this notebook,
we need to configure \texttt{matplotlib} (the plotting library) to work
within the notebook properly:

    Matplotlib documentation can be found
\href{https://matplotlib.org/contents.html}{\texttt{here}}.

    \begin{tcolorbox}[breakable, size=fbox, boxrule=1pt, pad at break*=1mm,colback=cellbackground, colframe=cellborder]
\prompt{In}{incolor}{7}{\boxspacing}
\begin{Verbatim}[commandchars=\\\{\}]
\PY{o}{\PYZpc{}}\PY{k}{matplotlib} inline
\PY{n}{plot} \PY{o}{=} \PY{n}{ldata}\PY{o}{.}\PY{n}{plot}\PY{p}{(}\PY{p}{)}
\end{Verbatim}
\end{tcolorbox}

    \begin{center}
    \adjustimage{max size={0.9\linewidth}{0.9\paperheight}}{output_18_0.png}
    \end{center}
    { \hspace*{\fill} \\}
    
    Notes: There are alternative ways to access the GWOSC data.

\begin{itemize}
\tightlist
\item
  \href{https://losc.ligo.org/s/sample_code/readligo.py}{\texttt{readligo}}
  is a light-weight Python module that returns the time series into a
  Numpy array.
\item
  The \href{http://github.com/ligo-cbc/pycbc}{PyCBC} package has the
  \texttt{pycbc.frame.query\_and\_read\_frame} and
  \texttt{pycbc.frame.read\_frame} methods. We use
  \href{http://github.com/ligo-cbc/pycbc}{PyCBC} in Tutorial 2.1, 2.2
  and 2.3.
\end{itemize}

    \hypertarget{handling-data-in-the-frequency-domain-using-the-fourier-transform}{%
\subsection{Handling data in the frequency domain using the Fourier
transform}\label{handling-data-in-the-frequency-domain-using-the-fourier-transform}}

The \href{https://en.wikipedia.org/wiki/Fourier_transform}{Fourier
transform} is a widely-used mathematical tool to expose the
frequency-domain content of a time-domain signal, meaning we can see
which frequencies contain lots of power, and which have less.

We can calculate the Fourier transform of our \texttt{TimeSeries} using
the
\href{https://gwpy.github.io/docs/stable/api/gwpy.timeseries.TimeSeries.html\#gwpy.timeseries.TimeSeries.fft}{\texttt{fft()}}
method:

    \begin{tcolorbox}[breakable, size=fbox, boxrule=1pt, pad at break*=1mm,colback=cellbackground, colframe=cellborder]
\prompt{In}{incolor}{8}{\boxspacing}
\begin{Verbatim}[commandchars=\\\{\}]
\PY{n}{fft} \PY{o}{=} \PY{n}{ldata}\PY{o}{.}\PY{n}{fft}\PY{p}{(}\PY{p}{)}
\PY{n+nb}{print}\PY{p}{(}\PY{n}{fft}\PY{p}{)}
\end{Verbatim}
\end{tcolorbox}

    \begin{Verbatim}[commandchars=\\\{\}]
FrequencySeries([-1.45894410e-21+0.00000000e+00j,
                 -2.91834926e-21-4.52905657e-23j,
                 -2.91973330e-21-9.06202944e-23j, {\ldots},
                 -2.38723982e-23+4.67871178e-26j,
                 -2.38345362e-23+1.80394058e-26j,
                 -2.38457175e-23+0.00000000e+00j]
                unit: dimensionless,
                f0: 0.0 Hz,
                df: 0.1 Hz,
                epoch: 1239082257.0,
                name: Strain,
                channel: None)
    \end{Verbatim}

    The result is a
\href{https://gwpy.github.io/docs/stable/frequencyseries/}{\texttt{FrequencySeries}},
with complex amplitude, representing the amplitude and phase of each
frequency in our data. We can use \texttt{abs()} to extract the
amplitude and plot that:

    \begin{tcolorbox}[breakable, size=fbox, boxrule=1pt, pad at break*=1mm,colback=cellbackground, colframe=cellborder]
\prompt{In}{incolor}{9}{\boxspacing}
\begin{Verbatim}[commandchars=\\\{\}]
\PY{n}{plot} \PY{o}{=} \PY{n}{fft}\PY{o}{.}\PY{n}{abs}\PY{p}{(}\PY{p}{)}\PY{o}{.}\PY{n}{plot}\PY{p}{(}\PY{n}{xscale}\PY{o}{=}\PY{l+s+s2}{\PYZdq{}}\PY{l+s+s2}{log}\PY{l+s+s2}{\PYZdq{}}\PY{p}{,} \PY{n}{yscale}\PY{o}{=}\PY{l+s+s2}{\PYZdq{}}\PY{l+s+s2}{log}\PY{l+s+s2}{\PYZdq{}}\PY{p}{)}
\PY{n}{plot}\PY{o}{.}\PY{n}{show}\PY{p}{(}\PY{n}{warn}\PY{o}{=}\PY{k+kc}{False}\PY{p}{)}
\end{Verbatim}
\end{tcolorbox}

    \begin{center}
    \adjustimage{max size={0.9\linewidth}{0.9\paperheight}}{output_23_0.png}
    \end{center}
    { \hspace*{\fill} \\}
    
    This doesn't look correct at all! The problem is that the FFT works
under the assumption that our data are periodic, which means that the
edges of our data look like discontinuities when transformed. We need to
apply a window function to our time-domain data before transforming,
which we can do using the
\href{https://docs.scipy.org/doc/scipy/reference/signal.html}{\texttt{scipy.signal}}
module:

    \begin{tcolorbox}[breakable, size=fbox, boxrule=1pt, pad at break*=1mm,colback=cellbackground, colframe=cellborder]
\prompt{In}{incolor}{10}{\boxspacing}
\begin{Verbatim}[commandchars=\\\{\}]
\PY{k+kn}{from} \PY{n+nn}{scipy}\PY{n+nn}{.}\PY{n+nn}{signal} \PY{k+kn}{import} \PY{n}{get\PYZus{}window}
\PY{n}{window} \PY{o}{=} \PY{n}{get\PYZus{}window}\PY{p}{(}\PY{l+s+s1}{\PYZsq{}}\PY{l+s+s1}{hann}\PY{l+s+s1}{\PYZsq{}}\PY{p}{,} \PY{n}{ldata}\PY{o}{.}\PY{n}{size}\PY{p}{)}
\PY{n}{lwin} \PY{o}{=} \PY{n}{ldata} \PY{o}{*} \PY{n}{window}
\end{Verbatim}
\end{tcolorbox}

    Let's try our transform again and see what we get

    \begin{tcolorbox}[breakable, size=fbox, boxrule=1pt, pad at break*=1mm,colback=cellbackground, colframe=cellborder]
\prompt{In}{incolor}{11}{\boxspacing}
\begin{Verbatim}[commandchars=\\\{\}]
\PY{n}{fftamp} \PY{o}{=} \PY{n}{lwin}\PY{o}{.}\PY{n}{fft}\PY{p}{(}\PY{p}{)}\PY{o}{.}\PY{n}{abs}\PY{p}{(}\PY{p}{)}
\PY{n}{plot} \PY{o}{=} \PY{n}{fftamp}\PY{o}{.}\PY{n}{plot}\PY{p}{(}\PY{n}{xscale}\PY{o}{=}\PY{l+s+s2}{\PYZdq{}}\PY{l+s+s2}{log}\PY{l+s+s2}{\PYZdq{}}\PY{p}{,} \PY{n}{yscale}\PY{o}{=}\PY{l+s+s2}{\PYZdq{}}\PY{l+s+s2}{log}\PY{l+s+s2}{\PYZdq{}}\PY{p}{)}
\PY{n}{plot}\PY{o}{.}\PY{n}{show}\PY{p}{(}\PY{n}{warn}\PY{o}{=}\PY{k+kc}{False}\PY{p}{)}
\end{Verbatim}
\end{tcolorbox}

    \begin{center}
    \adjustimage{max size={0.9\linewidth}{0.9\paperheight}}{output_27_0.png}
    \end{center}
    { \hspace*{\fill} \\}
    
    This looks a little more like what we expect for the amplitude spectral
density of a gravitational-wave detector.

    \hypertarget{calculating-the-power-spectral-density}{%
\subsection{Calculating the power spectral
density}\label{calculating-the-power-spectral-density}}

In practice, we typically use a large number of FFTs to estimate an
average power spectral density over a long period of data. We can do
this using the
\href{https://gwpy.github.io/docs/stable/api/gwpy.timeseries.TimeSeries.html\#gwpy.timeseries.TimeSeries.asd}{\texttt{asd()}}
method, which uses
\href{https://en.wikipedia.org/wiki/Welch\%27s_method}{Welch's method}
to combine FFTs of overlapping, windowed chunks of data.

    \begin{tcolorbox}[breakable, size=fbox, boxrule=1pt, pad at break*=1mm,colback=cellbackground, colframe=cellborder]
\prompt{In}{incolor}{12}{\boxspacing}
\begin{Verbatim}[commandchars=\\\{\}]
\PY{n}{asd} \PY{o}{=} \PY{n}{ldata}\PY{o}{.}\PY{n}{asd}\PY{p}{(}\PY{n}{fftlength}\PY{o}{=}\PY{l+m+mi}{4}\PY{p}{,} \PY{n}{method}\PY{o}{=}\PY{l+s+s2}{\PYZdq{}}\PY{l+s+s2}{median}\PY{l+s+s2}{\PYZdq{}}\PY{p}{)}
\PY{n}{plot} \PY{o}{=} \PY{n}{asd}\PY{o}{.}\PY{n}{plot}\PY{p}{(}\PY{p}{)}
\PY{n}{plot}\PY{o}{.}\PY{n}{show}\PY{p}{(}\PY{n}{warn}\PY{o}{=}\PY{k+kc}{False}\PY{p}{)}
\end{Verbatim}
\end{tcolorbox}

    \begin{center}
    \adjustimage{max size={0.9\linewidth}{0.9\paperheight}}{output_30_0.png}
    \end{center}
    { \hspace*{\fill} \\}
    
    \begin{tcolorbox}[breakable, size=fbox, boxrule=1pt, pad at break*=1mm,colback=cellbackground, colframe=cellborder]
\prompt{In}{incolor}{13}{\boxspacing}
\begin{Verbatim}[commandchars=\\\{\}]
\PY{n}{ax} \PY{o}{=} \PY{n}{plot}\PY{o}{.}\PY{n}{gca}\PY{p}{(}\PY{p}{)}
\PY{n}{ax}\PY{o}{.}\PY{n}{set\PYZus{}xlim}\PY{p}{(}\PY{l+m+mi}{10}\PY{p}{,} \PY{l+m+mi}{1400}\PY{p}{)}
\PY{n}{ax}\PY{o}{.}\PY{n}{set\PYZus{}ylim}\PY{p}{(}\PY{l+m+mf}{1e\PYZhy{}24}\PY{p}{,} \PY{l+m+mf}{1e\PYZhy{}20}\PY{p}{)}
\PY{n}{plot}
\end{Verbatim}
\end{tcolorbox}
 
            
\prompt{Out}{outcolor}{13}{}
    
    \begin{center}
    \adjustimage{max size={0.9\linewidth}{0.9\paperheight}}{output_31_0.png}
    \end{center}
    { \hspace*{\fill} \\}
    

    The ASD is a standard tool used to study the frequency-domain
sensitivity of a gravitational-wave detector. For the LIGO-Livingston
data we loaded, we can see large spikes at certain frequencies,
including

\begin{itemize}
\tightlist
\item
  \textasciitilde300 Hz
\item
  \textasciitilde500 Hz
\item
  \textasciitilde1000 Hz
\end{itemize}

The \href{https://www.gw-openscience.org/o2speclines/}{O2 spectral
lines} page on GWOSC describes a number of these spectral features for
O2, with some of them being forced upon us, and some being deliberately
introduced to help with interferometer control.

Loading more data allows for more FFTs to be averaged during the ASD
calculation, meaning random variations get averaged out, and we can see
more detail:

    \begin{tcolorbox}[breakable, size=fbox, boxrule=1pt, pad at break*=1mm,colback=cellbackground, colframe=cellborder]
\prompt{In}{incolor}{14}{\boxspacing}
\begin{Verbatim}[commandchars=\\\{\}]
\PY{n}{ldata2} \PY{o}{=} \PY{n}{TimeSeries}\PY{o}{.}\PY{n}{fetch\PYZus{}open\PYZus{}data}\PY{p}{(}\PY{l+s+s1}{\PYZsq{}}\PY{l+s+s1}{L1}\PY{l+s+s1}{\PYZsq{}}\PY{p}{,} \PY{n+nb}{int}\PY{p}{(}\PY{n}{gps}\PY{p}{)}\PY{o}{\PYZhy{}}\PY{l+m+mi}{512}\PY{p}{,} \PY{n+nb}{int}\PY{p}{(}\PY{n}{gps}\PY{p}{)}\PY{o}{+}\PY{l+m+mi}{512}\PY{p}{,} \PY{n}{cache}\PY{o}{=}\PY{k+kc}{True}\PY{p}{)}
\PY{n}{lasd2} \PY{o}{=} \PY{n}{ldata2}\PY{o}{.}\PY{n}{asd}\PY{p}{(}\PY{n}{fftlength}\PY{o}{=}\PY{l+m+mi}{4}\PY{p}{,} \PY{n}{method}\PY{o}{=}\PY{l+s+s2}{\PYZdq{}}\PY{l+s+s2}{median}\PY{l+s+s2}{\PYZdq{}}\PY{p}{)}
\PY{n}{plot} \PY{o}{=} \PY{n}{lasd2}\PY{o}{.}\PY{n}{plot}\PY{p}{(}\PY{p}{)}
\PY{n}{ax} \PY{o}{=} \PY{n}{plot}\PY{o}{.}\PY{n}{gca}\PY{p}{(}\PY{p}{)}
\PY{n}{ax}\PY{o}{.}\PY{n}{set\PYZus{}xlim}\PY{p}{(}\PY{l+m+mi}{10}\PY{p}{,} \PY{l+m+mi}{1400}\PY{p}{)}
\PY{n}{ax}\PY{o}{.}\PY{n}{set\PYZus{}ylim}\PY{p}{(}\PY{l+m+mf}{1e\PYZhy{}24}\PY{p}{,} \PY{l+m+mf}{1e\PYZhy{}20}\PY{p}{)}
\PY{n}{plot}\PY{o}{.}\PY{n}{show}\PY{p}{(}\PY{n}{warn}\PY{o}{=}\PY{k+kc}{False}\PY{p}{)}
\end{Verbatim}
\end{tcolorbox}

    \begin{center}
    \adjustimage{max size={0.9\linewidth}{0.9\paperheight}}{output_33_0.png}
    \end{center}
    { \hspace*{\fill} \\}
    
    Now we can see some more features, including sets of lines around
\textasciitilde30 Hz and \textasciitilde65 Hz, and some more isolated
lines through the more sensitive region.

For comparison, we can load the LIGO-Hanford data and plot that as well:

    \begin{tcolorbox}[breakable, size=fbox, boxrule=1pt, pad at break*=1mm,colback=cellbackground, colframe=cellborder]
\prompt{In}{incolor}{15}{\boxspacing}
\begin{Verbatim}[commandchars=\\\{\}]
\PY{c+c1}{\PYZsh{} get Hanford data}
\PY{n}{hdata2} \PY{o}{=} \PY{n}{TimeSeries}\PY{o}{.}\PY{n}{fetch\PYZus{}open\PYZus{}data}\PY{p}{(}\PY{l+s+s1}{\PYZsq{}}\PY{l+s+s1}{H1}\PY{l+s+s1}{\PYZsq{}}\PY{p}{,} \PY{n+nb}{int}\PY{p}{(}\PY{n}{gps}\PY{p}{)}\PY{o}{\PYZhy{}}\PY{l+m+mi}{512}\PY{p}{,} \PY{n+nb}{int}\PY{p}{(}\PY{n}{gps}\PY{p}{)}\PY{o}{+}\PY{l+m+mi}{512}\PY{p}{,} \PY{n}{cache}\PY{o}{=}\PY{k+kc}{True}\PY{p}{)}
\PY{n}{hasd2} \PY{o}{=} \PY{n}{hdata2}\PY{o}{.}\PY{n}{asd}\PY{p}{(}\PY{n}{fftlength}\PY{o}{=}\PY{l+m+mi}{4}\PY{p}{,} \PY{n}{method}\PY{o}{=}\PY{l+s+s2}{\PYZdq{}}\PY{l+s+s2}{median}\PY{l+s+s2}{\PYZdq{}}\PY{p}{)}

\PY{c+c1}{\PYZsh{} get Virgo data}
\PY{n}{vdata2} \PY{o}{=} \PY{n}{TimeSeries}\PY{o}{.}\PY{n}{fetch\PYZus{}open\PYZus{}data}\PY{p}{(}\PY{l+s+s1}{\PYZsq{}}\PY{l+s+s1}{V1}\PY{l+s+s1}{\PYZsq{}}\PY{p}{,} \PY{n+nb}{int}\PY{p}{(}\PY{n}{gps}\PY{p}{)}\PY{o}{\PYZhy{}}\PY{l+m+mi}{512}\PY{p}{,} \PY{n+nb}{int}\PY{p}{(}\PY{n}{gps}\PY{p}{)}\PY{o}{+}\PY{l+m+mi}{512}\PY{p}{,} \PY{n}{cache}\PY{o}{=}\PY{k+kc}{True}\PY{p}{)}
\PY{n}{vasd2} \PY{o}{=} \PY{n}{vdata2}\PY{o}{.}\PY{n}{asd}\PY{p}{(}\PY{n}{fftlength}\PY{o}{=}\PY{l+m+mi}{4}\PY{p}{,} \PY{n}{method}\PY{o}{=}\PY{l+s+s2}{\PYZdq{}}\PY{l+s+s2}{median}\PY{l+s+s2}{\PYZdq{}}\PY{p}{)}

\PY{c+c1}{\PYZsh{} and plot using standard colours}
\PY{n}{ax}\PY{o}{.}\PY{n}{plot}\PY{p}{(}\PY{n}{hasd2}\PY{p}{,} \PY{n}{label}\PY{o}{=}\PY{l+s+s1}{\PYZsq{}}\PY{l+s+s1}{LIGO\PYZhy{}Hanford}\PY{l+s+s1}{\PYZsq{}}\PY{p}{,} \PY{n}{color}\PY{o}{=}\PY{l+s+s1}{\PYZsq{}}\PY{l+s+s1}{gwpy:ligo\PYZhy{}hanford}\PY{l+s+s1}{\PYZsq{}}\PY{p}{)}
\PY{n}{ax}\PY{o}{.}\PY{n}{plot}\PY{p}{(}\PY{n}{vasd2}\PY{p}{,} \PY{n}{label}\PY{o}{=}\PY{l+s+s1}{\PYZsq{}}\PY{l+s+s1}{Virgo}\PY{l+s+s1}{\PYZsq{}}\PY{p}{,} \PY{n}{color}\PY{o}{=}\PY{l+s+s1}{\PYZsq{}}\PY{l+s+s1}{gwpy:virgo}\PY{l+s+s1}{\PYZsq{}}\PY{p}{)}

\PY{c+c1}{\PYZsh{} update the Livingston line to use standard colour, and have a label}
\PY{n}{lline} \PY{o}{=} \PY{n}{ax}\PY{o}{.}\PY{n}{lines}\PY{p}{[}\PY{l+m+mi}{0}\PY{p}{]}
\PY{n}{lline}\PY{o}{.}\PY{n}{set\PYZus{}color}\PY{p}{(}\PY{l+s+s1}{\PYZsq{}}\PY{l+s+s1}{gwpy:ligo\PYZhy{}livingston}\PY{l+s+s1}{\PYZsq{}}\PY{p}{)}  \PY{c+c1}{\PYZsh{} change colour of Livingston data}
\PY{n}{lline}\PY{o}{.}\PY{n}{set\PYZus{}label}\PY{p}{(}\PY{l+s+s1}{\PYZsq{}}\PY{l+s+s1}{LIGO\PYZhy{}Livingston}\PY{l+s+s1}{\PYZsq{}}\PY{p}{)}

\PY{n}{ax}\PY{o}{.}\PY{n}{set\PYZus{}ylabel}\PY{p}{(}\PY{l+s+sa}{r}\PY{l+s+s1}{\PYZsq{}}\PY{l+s+s1}{Strain noise [\PYZdl{}1/}\PY{l+s+s1}{\PYZbs{}}\PY{l+s+s1}{sqrt}\PY{l+s+s1}{\PYZob{}}\PY{l+s+s1}{\PYZbs{}}\PY{l+s+s1}{mathrm}\PY{l+s+si}{\PYZob{}Hz\PYZcb{}}\PY{l+s+s1}{\PYZcb{}\PYZdl{}]}\PY{l+s+s1}{\PYZsq{}}\PY{p}{)}
\PY{n}{ax}\PY{o}{.}\PY{n}{legend}\PY{p}{(}\PY{p}{)}
\PY{n}{plot}
\end{Verbatim}
\end{tcolorbox}
 
            
\prompt{Out}{outcolor}{15}{}
    
    \begin{center}
    \adjustimage{max size={0.9\linewidth}{0.9\paperheight}}{output_35_0.png}
    \end{center}
    { \hspace*{\fill} \\}
    

    Now we can see clearly the relative sensitivity of each LIGO instrument,
the common features between both, and those unique to each observatory.

    \hypertarget{challenges}{%
\section{Challenges:}\label{challenges}}

    \hypertarget{quiz-question-1}{%
\section{Quiz Question 1:}\label{quiz-question-1}}

La amplitud máxima en los datos de LIGO-Livingston ocurre
aproximadamente a los 5 segundos en la gráfica anterior y es imposible
de detectar a simple vista por encima del ruido de fondo. Grafique los
datos para el detector LIGO-Hanford alrededor del evento GW190412. Al
observar su nueva gráfica de LIGO-Hanford, ¿puede su ojo identificar un
pico de señal?

    \begin{tcolorbox}[breakable, size=fbox, boxrule=1pt, pad at break*=1mm,colback=cellbackground, colframe=cellborder]
\prompt{In}{incolor}{16}{\boxspacing}
\begin{Verbatim}[commandchars=\\\{\}]
\PY{n}{ldata2} \PY{o}{=} \PY{n}{TimeSeries}\PY{o}{.}\PY{n}{fetch\PYZus{}open\PYZus{}data}\PY{p}{(}\PY{l+s+s1}{\PYZsq{}}\PY{l+s+s1}{H1}\PY{l+s+s1}{\PYZsq{}}\PY{p}{,} \PY{n+nb}{int}\PY{p}{(}\PY{n}{gps}\PY{p}{)}\PY{o}{\PYZhy{}}\PY{l+m+mi}{512}\PY{p}{,} \PY{n+nb}{int}\PY{p}{(}\PY{n}{gps}\PY{p}{)}\PY{o}{+}\PY{l+m+mi}{512}\PY{p}{,} \PY{n}{cache}\PY{o}{=}\PY{k+kc}{True}\PY{p}{)}
\PY{n}{lasd2} \PY{o}{=} \PY{n}{ldata2}\PY{o}{.}\PY{n}{asd}\PY{p}{(}\PY{n}{fftlength}\PY{o}{=}\PY{l+m+mi}{4}\PY{p}{,} \PY{n}{method}\PY{o}{=}\PY{l+s+s2}{\PYZdq{}}\PY{l+s+s2}{median}\PY{l+s+s2}{\PYZdq{}}\PY{p}{)}
\PY{n}{plot} \PY{o}{=} \PY{n}{lasd2}\PY{o}{.}\PY{n}{plot}\PY{p}{(}\PY{p}{)}
\PY{n}{ax} \PY{o}{=} \PY{n}{plot}\PY{o}{.}\PY{n}{gca}\PY{p}{(}\PY{p}{)}
\PY{n}{ax}\PY{o}{.}\PY{n}{set\PYZus{}xlim}\PY{p}{(}\PY{l+m+mi}{10}\PY{p}{,} \PY{l+m+mi}{1400}\PY{p}{)}
\PY{n}{ax}\PY{o}{.}\PY{n}{set\PYZus{}ylim}\PY{p}{(}\PY{l+m+mf}{1e\PYZhy{}24}\PY{p}{,} \PY{l+m+mf}{1e\PYZhy{}20}\PY{p}{)}
\PY{n}{plot}\PY{o}{.}\PY{n}{show}\PY{p}{(}\PY{n}{warn}\PY{o}{=}\PY{k+kc}{False}\PY{p}{)}
\end{Verbatim}
\end{tcolorbox}

    \begin{center}
    \adjustimage{max size={0.9\linewidth}{0.9\paperheight}}{output_39_0.png}
    \end{center}
    { \hspace*{\fill} \\}
    
    No se puede detectar un pico de señal ya que la señal de ruido sigue
siendo muy grande para detectar las señales a simple vista

    \hypertarget{quiz-question-2}{%
\section{Quiz Question 2 :}\label{quiz-question-2}}

Cree una ASD (densidad espectral de amplitud) alrededor del momento de
un evento de O1, GW150914, para el detector L1. Compare esto con las ASD
alrededor de GW190412 para el detector L1. ¿Qué datos tienen menor ruido
---y por lo tanto son más sensibles--- alrededor de los 100 Hz?

    \begin{tcolorbox}[breakable, size=fbox, boxrule=1pt, pad at break*=1mm,colback=cellbackground, colframe=cellborder]
\prompt{In}{incolor}{17}{\boxspacing}
\begin{Verbatim}[commandchars=\\\{\}]
\PY{k+kn}{from} \PY{n+nn}{gwosc}\PY{n+nn}{.}\PY{n+nn}{datasets} \PY{k+kn}{import} \PY{n}{event\PYZus{}gps}
\PY{n}{gps} \PY{o}{=} \PY{n}{event\PYZus{}gps}\PY{p}{(}\PY{l+s+s1}{\PYZsq{}}\PY{l+s+s1}{GW150914}\PY{l+s+s1}{\PYZsq{}}\PY{p}{)}
\PY{n+nb}{print}\PY{p}{(}\PY{n}{gps}\PY{p}{)}
\end{Verbatim}
\end{tcolorbox}

    \begin{Verbatim}[commandchars=\\\{\}]
1126259462.4
    \end{Verbatim}

    \begin{tcolorbox}[breakable, size=fbox, boxrule=1pt, pad at break*=1mm,colback=cellbackground, colframe=cellborder]
\prompt{In}{incolor}{18}{\boxspacing}
\begin{Verbatim}[commandchars=\\\{\}]
\PY{n}{segment} \PY{o}{=} \PY{p}{(}\PY{n+nb}{int}\PY{p}{(}\PY{n}{gps}\PY{p}{)}\PY{o}{\PYZhy{}}\PY{l+m+mi}{5}\PY{p}{,} \PY{n+nb}{int}\PY{p}{(}\PY{n}{gps}\PY{p}{)}\PY{o}{+}\PY{l+m+mi}{5}\PY{p}{)}
\PY{n+nb}{print}\PY{p}{(}\PY{n}{segment}\PY{p}{)}
\end{Verbatim}
\end{tcolorbox}

    \begin{Verbatim}[commandchars=\\\{\}]
(1126259457, 1126259467)
    \end{Verbatim}

    \begin{tcolorbox}[breakable, size=fbox, boxrule=1pt, pad at break*=1mm,colback=cellbackground, colframe=cellborder]
\prompt{In}{incolor}{19}{\boxspacing}
\begin{Verbatim}[commandchars=\\\{\}]
\PY{k+kn}{from} \PY{n+nn}{gwpy}\PY{n+nn}{.}\PY{n+nn}{timeseries} \PY{k+kn}{import} \PY{n}{TimeSeries}
\PY{n}{ldata} \PY{o}{=} \PY{n}{TimeSeries}\PY{o}{.}\PY{n}{fetch\PYZus{}open\PYZus{}data}\PY{p}{(}\PY{l+s+s1}{\PYZsq{}}\PY{l+s+s1}{L1}\PY{l+s+s1}{\PYZsq{}}\PY{p}{,} \PY{o}{*}\PY{n}{segment}\PY{p}{,} \PY{n}{verbose}\PY{o}{=}\PY{k+kc}{True}\PY{p}{)}
\PY{n+nb}{print}\PY{p}{(}\PY{n}{ldata}\PY{p}{)}
\end{Verbatim}
\end{tcolorbox}

    \begin{Verbatim}[commandchars=\\\{\}]
Fetched 1 URLs from gwosc.org for [1126259457 .. 1126259467))
Reading data{\ldots} [Done]
TimeSeries([-9.31087178e-19, -9.75697062e-19, -1.01074940e-18,
            {\ldots}, -1.13460681e-18, -1.11414494e-18,
            -1.15931435e-18]
           unit: dimensionless,
           t0: 1126259457.0 s,
           dt: 0.000244140625 s,
           name: Strain,
           channel: None)
    \end{Verbatim}

    \begin{tcolorbox}[breakable, size=fbox, boxrule=1pt, pad at break*=1mm,colback=cellbackground, colframe=cellborder]
\prompt{In}{incolor}{20}{\boxspacing}
\begin{Verbatim}[commandchars=\\\{\}]
\PY{o}{\PYZpc{}}\PY{k}{matplotlib} inline
\PY{n}{plot} \PY{o}{=} \PY{n}{ldata}\PY{o}{.}\PY{n}{plot}\PY{p}{(}\PY{p}{)}
\end{Verbatim}
\end{tcolorbox}

    \begin{center}
    \adjustimage{max size={0.9\linewidth}{0.9\paperheight}}{output_45_0.png}
    \end{center}
    { \hspace*{\fill} \\}
    
    \begin{tcolorbox}[breakable, size=fbox, boxrule=1pt, pad at break*=1mm,colback=cellbackground, colframe=cellborder]
\prompt{In}{incolor}{21}{\boxspacing}
\begin{Verbatim}[commandchars=\\\{\}]
\PY{n}{fft} \PY{o}{=} \PY{n}{ldata}\PY{o}{.}\PY{n}{fft}\PY{p}{(}\PY{p}{)}
\PY{n+nb}{print}\PY{p}{(}\PY{n}{fft}\PY{p}{)}
\end{Verbatim}
\end{tcolorbox}

    \begin{Verbatim}[commandchars=\\\{\}]
FrequencySeries([-1.05219333e-18+0.00000000e+00j,
                 -9.36314427e-22+6.61514979e-23j,
                 -8.74693847e-22-3.26717194e-23j, {\ldots},
                  6.63325463e-24-1.10870240e-26j,
                  6.73663868e-24-7.82178218e-26j,
                  6.69598891e-24+0.00000000e+00j]
                unit: dimensionless,
                f0: 0.0 Hz,
                df: 0.1 Hz,
                epoch: 1126259457.0,
                name: Strain,
                channel: None)
    \end{Verbatim}

    \begin{tcolorbox}[breakable, size=fbox, boxrule=1pt, pad at break*=1mm,colback=cellbackground, colframe=cellborder]
\prompt{In}{incolor}{23}{\boxspacing}
\begin{Verbatim}[commandchars=\\\{\}]
\PY{n}{plot} \PY{o}{=} \PY{n}{fft}\PY{o}{.}\PY{n}{abs}\PY{p}{(}\PY{p}{)}\PY{o}{.}\PY{n}{plot}\PY{p}{(}\PY{n}{xscale}\PY{o}{=}\PY{l+s+s2}{\PYZdq{}}\PY{l+s+s2}{log}\PY{l+s+s2}{\PYZdq{}}\PY{p}{,} \PY{n}{yscale}\PY{o}{=}\PY{l+s+s2}{\PYZdq{}}\PY{l+s+s2}{log}\PY{l+s+s2}{\PYZdq{}}\PY{p}{)}
\PY{n}{plot}\PY{o}{.}\PY{n}{show}\PY{p}{(}\PY{n}{warn}\PY{o}{=}\PY{k+kc}{False}\PY{p}{)}
\end{Verbatim}
\end{tcolorbox}

    \begin{center}
    \adjustimage{max size={0.9\linewidth}{0.9\paperheight}}{output_47_0.png}
    \end{center}
    { \hspace*{\fill} \\}
    
    \begin{tcolorbox}[breakable, size=fbox, boxrule=1pt, pad at break*=1mm,colback=cellbackground, colframe=cellborder]
\prompt{In}{incolor}{24}{\boxspacing}
\begin{Verbatim}[commandchars=\\\{\}]
\PY{k+kn}{from} \PY{n+nn}{scipy}\PY{n+nn}{.}\PY{n+nn}{signal} \PY{k+kn}{import} \PY{n}{get\PYZus{}window}
\PY{n}{window} \PY{o}{=} \PY{n}{get\PYZus{}window}\PY{p}{(}\PY{l+s+s1}{\PYZsq{}}\PY{l+s+s1}{hann}\PY{l+s+s1}{\PYZsq{}}\PY{p}{,} \PY{n}{ldata}\PY{o}{.}\PY{n}{size}\PY{p}{)}
\PY{n}{lwin} \PY{o}{=} \PY{n}{ldata} \PY{o}{*} \PY{n}{window}
\end{Verbatim}
\end{tcolorbox}

    \begin{tcolorbox}[breakable, size=fbox, boxrule=1pt, pad at break*=1mm,colback=cellbackground, colframe=cellborder]
\prompt{In}{incolor}{25}{\boxspacing}
\begin{Verbatim}[commandchars=\\\{\}]
\PY{n}{fftamp} \PY{o}{=} \PY{n}{lwin}\PY{o}{.}\PY{n}{fft}\PY{p}{(}\PY{p}{)}\PY{o}{.}\PY{n}{abs}\PY{p}{(}\PY{p}{)}
\PY{n}{plot} \PY{o}{=} \PY{n}{fftamp}\PY{o}{.}\PY{n}{plot}\PY{p}{(}\PY{n}{xscale}\PY{o}{=}\PY{l+s+s2}{\PYZdq{}}\PY{l+s+s2}{log}\PY{l+s+s2}{\PYZdq{}}\PY{p}{,} \PY{n}{yscale}\PY{o}{=}\PY{l+s+s2}{\PYZdq{}}\PY{l+s+s2}{log}\PY{l+s+s2}{\PYZdq{}}\PY{p}{)}
\PY{n}{plot}\PY{o}{.}\PY{n}{show}\PY{p}{(}\PY{n}{warn}\PY{o}{=}\PY{k+kc}{False}\PY{p}{)}
\end{Verbatim}
\end{tcolorbox}

    \begin{center}
    \adjustimage{max size={0.9\linewidth}{0.9\paperheight}}{output_49_0.png}
    \end{center}
    { \hspace*{\fill} \\}
    
    \begin{tcolorbox}[breakable, size=fbox, boxrule=1pt, pad at break*=1mm,colback=cellbackground, colframe=cellborder]
\prompt{In}{incolor}{32}{\boxspacing}
\begin{Verbatim}[commandchars=\\\{\}]
\PY{n}{ax} \PY{o}{=} \PY{n}{plot}\PY{o}{.}\PY{n}{gca}\PY{p}{(}\PY{p}{)}
\PY{n}{ax}\PY{o}{.}\PY{n}{set\PYZus{}xlim}\PY{p}{(}\PY{l+m+mi}{10}\PY{p}{,} \PY{l+m+mi}{1400}\PY{p}{)}
\PY{n}{ax}\PY{o}{.}\PY{n}{set\PYZus{}ylim}\PY{p}{(}\PY{l+m+mf}{1e\PYZhy{}26}\PY{p}{,} \PY{l+m+mf}{1e\PYZhy{}20}\PY{p}{)}
\PY{n}{plot}
\end{Verbatim}
\end{tcolorbox}
 
            
\prompt{Out}{outcolor}{32}{}
    
    \begin{center}
    \adjustimage{max size={0.9\linewidth}{0.9\paperheight}}{output_50_0.png}
    \end{center}
    { \hspace*{\fill} \\}
    

    \begin{tcolorbox}[breakable, size=fbox, boxrule=1pt, pad at break*=1mm,colback=cellbackground, colframe=cellborder]
\prompt{In}{incolor}{35}{\boxspacing}
\begin{Verbatim}[commandchars=\\\{\}]
\PY{n}{ldata2} \PY{o}{=} \PY{n}{TimeSeries}\PY{o}{.}\PY{n}{fetch\PYZus{}open\PYZus{}data}\PY{p}{(}\PY{l+s+s1}{\PYZsq{}}\PY{l+s+s1}{L1}\PY{l+s+s1}{\PYZsq{}}\PY{p}{,} \PY{n+nb}{int}\PY{p}{(}\PY{n}{gps}\PY{p}{)}\PY{o}{\PYZhy{}}\PY{l+m+mi}{512}\PY{p}{,} \PY{n+nb}{int}\PY{p}{(}\PY{n}{gps}\PY{p}{)}\PY{o}{+}\PY{l+m+mi}{512}\PY{p}{,} \PY{n}{cache}\PY{o}{=}\PY{k+kc}{True}\PY{p}{)}
\PY{n}{lasd2} \PY{o}{=} \PY{n}{ldata2}\PY{o}{.}\PY{n}{asd}\PY{p}{(}\PY{n}{fftlength}\PY{o}{=}\PY{l+m+mi}{4}\PY{p}{,} \PY{n}{method}\PY{o}{=}\PY{l+s+s2}{\PYZdq{}}\PY{l+s+s2}{median}\PY{l+s+s2}{\PYZdq{}}\PY{p}{)}
\PY{n}{plot} \PY{o}{=} \PY{n}{lasd2}\PY{o}{.}\PY{n}{plot}\PY{p}{(}\PY{p}{)}
\PY{n}{ax} \PY{o}{=} \PY{n}{plot}\PY{o}{.}\PY{n}{gca}\PY{p}{(}\PY{p}{)}
\PY{n}{ax}\PY{o}{.}\PY{n}{set\PYZus{}xlim}\PY{p}{(}\PY{l+m+mi}{10}\PY{p}{,} \PY{l+m+mi}{1400}\PY{p}{)}
\PY{n}{ax}\PY{o}{.}\PY{n}{set\PYZus{}ylim}\PY{p}{(}\PY{l+m+mf}{1e\PYZhy{}24}\PY{p}{,} \PY{l+m+mf}{1e\PYZhy{}20}\PY{p}{)}
\PY{n}{plot}\PY{o}{.}\PY{n}{show}\PY{p}{(}\PY{n}{warn}\PY{o}{=}\PY{k+kc}{False}\PY{p}{)}
\end{Verbatim}
\end{tcolorbox}

    \begin{center}
    \adjustimage{max size={0.9\linewidth}{0.9\paperheight}}{output_51_0.png}
    \end{center}
    { \hspace*{\fill} \\}
    
    Al rededor de los 100 Hz el evento GW190412 presenta menor ruido que el
GW150914


    % Add a bibliography block to the postdoc
    
    
    
\end{document}
